% !PRINTED 167
\weberpage{167}

%<<CONT
divisor occur, only that here the indecomposable functions
take over the role of the primes.

We add one more theorem of this kind:

\gesperrt{IV. If $u$, $v$ are relatively prime and $Uu$, $Uv$ are
divisible by $w$, then $U$ is divisible by $w$.}

For if one decomposes the integral rational functions $U, u, v, w$
into their indecomposable factors, then any factor of
$w$, since it cannot occur in $u$ and $v$ at the same time, must
divide $U$. If one removes it from $U$ and $w$, one can conclude
in the same way for a next factor of $w$, and so on.

We now consider integral rational functions of one vari\-
able $t$
\[
f(t) = u_0 t^m + u_1 t^{m-1} + \cdots + u_{m-1} t + u_m,
\]
whose coefficients are integral rational functions of $n$ vari\-
ables $x$, of which $t$ is independent.

If the coefficients $u_0, u_1 \ldots u_m$ have no common
divisor, then $f(t)$ is called \gesperrt{primitive}; in the other case the
greatest common divisor of the coefficients $u_0, u_1 \ldots u_m$
is called \gesperrt{the divisor of the function} $f(t)$ (cf. §.\,2).

We conclude, always under the assumption of the validity
of I:

\gesperrt{V. The product of two primitive functions $f(t)$
and $\varphi(t)$ is again a primitive function, and
the divisor of a product of two imprimitive
functions is equal to the product of the divisors
of the two factors.}

The proof is exactly as for the corresponding theorem
in §.\,2.

Let
\weq{1}{
\begin{aligned}
f(t) &= u_0 t^m + u_1 t^{m-1} + \cdots + u_m \\
\varphi(t) &= v_0 t^\mu + v_1 t^{\mu-1} + \cdots + v_\mu
\end{aligned}
}
be two primitive functions and
\weq{2}{
F(t) = U_0 t^{m+\mu} + U_1 t^{m+\mu-1} + \cdots + U_{m+\mu}
}
their product. Then, if $r$ and $s$ denote any two numbers from
the series $0, 1 \ldots m$, and $0, 1, 2 \ldots \mu$ (§.\,2), we have
\weq{3}{
\begin{aligned}
U_{r+s} &= u_r v_s + u_{r-1} v_{s+1} + \cdots \\
&\qquad + u_{r+1} v_{s-1} + \cdots
\end{aligned}
}
%CONT>>
